\par\bigskip
\begin{center}
\rule{.88\textwidth}{.3pt}\par
Göttingische gelehrte Anzeigen.\quad 1823 Februar 24.\par
\rule{.88\textwidth}{.3pt}
\end{center}
\pdfbookmark[1]{1823 Februar 24}{anzeige-1823-02-24}

Eine am 2. Febr. der Königl. Societät vom Hrn. Hofr. \textsc{Gauss} überreichte Vorlesung, überschrieben
\begin{center}
\textit{Theoria combinationis observationum erroribus minimis obnoxiae, pars posterior,}
\end{center}
\noindent steht im unmittelbaren Zusammenhange mit einer frühern, wovon in diesen Blättern [1821 Februar 26] eine Anzeige gegeben ist. Wir bringen darüber nur kurz in Erinnerung, dass ihr Zweck war, die sogenannte Methode der kleinsten Quadrate auf eine neue Art zu begründen, wobei diese Methode nicht näherungsweise, sondern in mathematischer Schärfe, nicht mit der Beschränkung auf den Fall einer sehr grossen Anzahl von Beobachtungen, und nicht abhängig von einem hypothetischen Gesetze für die Wahrscheinlichkeit der Beobachtungsfehler, sondern in vollkommener Allgemeinheit, als die zweckmässigste Combinationsart der Beobachtungen erscheint. Der gegenwärtige zweite Theil der Untersuchung enthält nun eine weitere Ausführung dieser Lehre in einer Reihe von Lehrsätzen und Problemen, die damit in genauester Verbindung stehen. Es würde der Einrich%
